% !TeX program = lualatex
% Compile twice: lualatex 84.tex
% Standalone research notebook; original section preserved.
% Research additions: 27 September 2026. No external TeX input or BibTeX file.
\documentclass[11pt,a4paper]{article}
\usepackage[margin=24mm,headheight=14pt]{geometry}
\usepackage{fontspec}
\setmainfont{Latin Modern Roman}
\setsansfont{Latin Modern Sans}
\setmonofont{Latin Modern Mono}
\usepackage{amsmath,amssymb,mathtools}
\usepackage{unicode-math}
\setmathfont{Latin Modern Math}
\usepackage{microtype}
\usepackage{enumitem,xurl,needspace}
\usepackage[dvipsnames]{xcolor}
\usepackage{hyperref}
\hypersetup{unicode=true,colorlinks=true,linkcolor=MidnightBlue,urlcolor=MidnightBlue,
 pdftitle={Research notebook - Problem 84},pdfauthor={Open Problems Math},bookmarksnumbered=true}
\usepackage{bookmark}
\usepackage{titlesec}
\titleformat{\section}{\Large\bfseries\raggedright}{\thesection}{0.7em}{}
\usepackage{fancyhdr}
\pagestyle{fancy}\fancyhf{}
\fancyhead[L]{\small\sffamily Open Problems Math | Research notebook}
\fancyhead[R]{\small\sffamily Problem 84 | 27 September 2026}
\fancyfoot[C]{\thepage}
\setlength{\parindent}{0pt}
\setlength{\parskip}{5pt plus 1pt}
\setlength{\emergencystretch}{3em}
\setcounter{tocdepth}{1}
\setcounter{secnumdepth}{1}
\setlist[itemize]{leftmargin=3em,itemsep=5pt,topsep=4pt,parsep=0pt}
\allowdisplaybreaks
\numberwithin{equation}{section}
\newcommand{\N}{\mathbb{N}}
\newcommand{\Z}{\mathbb{Z}}
\newcommand{\Q}{\mathbb{Q}}
\newcommand{\R}{\mathbb{R}}
\newcommand{\C}{\mathbb{C}}
\newcommand{\F}{\mathbb{F}}
\newcommand{\A}{\mathbb{A}}
\newcommand{\PP}{\mathbb P}
\newcommand{\E}{\mathbb{E}}
\newcommand{\eps}{\varepsilon}
\newcommand{\dd}{\,\mathrm d}
\newcommand{\sref}[2]{\hyperlink{source:#1:#2}{[S#2]}}
\newcommand{\eref}[2]{\hyperlink{extra:#1:#2}{[E#2]}}
\newif\ifshowcatalogue
\showcataloguetrue
\newcommand{\cataloguescope}[1]{\ifshowcatalogue
 \paragraph{Exact scope in the supplied catalogue.}
 \begin{quote}\small #1\end{quote}\fi}
\begin{document}
\setcounter{section}{83}
\section{\texorpdfstring{Sarnak's Mobius randomness conjecture for zero-entropy systems}{Sarnak's Mobius randomness conjecture for zero-entropy systems}}\label{84}
\subsection{Definitions and mathematical statement}
The Möbius function is $\mu(n)=0$ when a prime square divides $n$, and $\mu(n)=(-1)^r$ when $n$ is a product of $r$ distinct primes. A compact metric system consists of a continuous map $T:X\to X$; zero topological entropy means its orbit complexity has zero exponential growth rate. The assertion is
\[
 \frac1N\sum_{n=1}^{N}\mu(n)f(T^nx)\longrightarrow0
 \qquad(f\in C(X),\ x\in X).
\]
Every point is quantified, not just almost every point under a chosen invariant probability measure. The test functions may be complex-valued.
\par\smallskip\textit{Formulation and concept references:} \sref{84}{1}.
\cataloguescope{Prove that for every compact metric topological dynamical system (X,T) of zero topological entropy, every continuous f on X, and every x in X, the average N\textasciicircum{}(-1) sum\_\{n<=N\} mu(n) f(T\textasciicircum{}n x) tends to zero as N tends to infinity.}
\subsection{Short English statement}
Is the Möbius function asymptotically uncorrelated with every continuous observable of every zero-entropy dynamical system?
% BEGIN RESEARCH_ADDITION ATTEMPT_84
\subsection{Research attempt}

\paragraph{Current frontier and scope.}
The 2026 infinite-dimensional torus paper addresses its specified torus
class \sref{84}{1}, not every zero-entropy system. Frantzikinakis--Host
prove a logarithmic version for a broad class including uniquely ergodic
zero-entropy systems; its weighting and hypotheses cannot be discarded
\eref{84}{1}. The familiar Chowla-to-Sarnak implication motivates the
attempt, but here its ordinary mixed-power arithmetic assumption is
written explicitly, rather than confused with an already established
logarithmic theorem or the exact Liouville statement of Section~77.

\paragraph{Attack route.}
Topological zero entropy gives subexponentially many orbit blocks, suggesting
a covering argument. A second-moment estimate alone can be overwhelmed by
that many blocks. The route pursued combines the covering with arbitrarily
high fixed moments, choosing the moment order after the block length but
before the averaging limit. This produces an explicit entropy--moment
tradeoff and identifies a sufficient arithmetic input. A sparse-set
alignment proposal is also checked against the density normalization.

\paragraph{Attempt: reducing every orbit to finitely many blocks.}
Fix a continuous complex-valued $f$ and scale so that $\|f\|_\infty\leq1$;
the zero function is immediate. For a block length $L$ and $\varepsilon>0$,
choose a finite set $\mathcal V_{L,\varepsilon}\subset\C^L$ of actual
observable blocks such that every
\[
 (f(Ty),\ldots,f(T^Ly)),\qquad y\in X,
\]
is within $\varepsilon$ in supremum norm of a member. Put
$M_{L,\varepsilon}=|\mathcal V_{L,\varepsilon}|$.
Uniform continuity of $f$ and zero topological entropy give
\begin{equation}\label{eq84:entropy}
 \log M_{L,\varepsilon}=o(L)\quad(L\to\infty)
\end{equation}
for every fixed $\varepsilon$. This is a covering-number statement for all
points of the compact space, not an almost-everywhere orbit theorem.

For every $x$, sliding the length-$L$ window gives
\begin{equation}\label{eq84:cover}
 \left|\frac1N\sum_{n=1}^N\mu(n)f(T^nx)\right|
 \leq\varepsilon+
 \frac1N\sum_{m=0}^{N-1}\max_{v\in\mathcal V_{L,\varepsilon}}
       \left|\frac1L\sum_{j=1}^L\mu(m+j)v_j\right|
       +O(L/N).
\end{equation}
To check the boundary error, average the original products over all windows.
Each interior index is counted exactly $L$ times, and the two ends together
contain at most $2L$ indices with smaller or extra weight. Their summands
have absolute value at most one. Replacing each orbit block by its net
representative changes its normalized sum by at most $\varepsilon$.
This proves \eqref{eq84:cover} with a bound independent of the chosen $x$.

\paragraph{Attempt: an explicit arithmetic hypothesis and moment bound.}
Assume the following \emph{ordinary mixed-power M\"obius Chowla statement}:
for every finite collection of distinct nonnegative shifts $h_i$ and
exponents $e_i\in\{1,2\}$, with at least one $e_i=1$,
\begin{equation}\label{eq84:chowla}
 \frac1N\sum_{m=1}^N\prod_i\mu(m+h_i)^{e_i}\longrightarrow0.
\end{equation}
This is an additional unresolved hypothesis, not an established consequence
of the logarithmic results. Its role in the classical implication is
consistent with the averaging distinctions in \eref{84}{1} and
\eref{84}{2}.

For fixed $L$, integer $r\geq1$, and $|v_j|\leq1$, expand the $2r$th
power of $|\sum_{j=1}^L\mu(m+j)v_j|$. Any term in which some index appears
an odd number of times has mean tending to zero by \eqref{eq84:chowla},
since odd positive powers of $\mu$ equal $\mu$ and even positive powers
equal $\mu^2$. A surviving term must have all multiplicities even, and
its coefficient and summand each have absolute value at most one.
The number of ordered $2r$-tuples with all multiplicities even is at most
\[
 (2r-1)!!\,L^r.
\]
Indeed, pair the $2r$ positions in $(2r-1)!!$ ways and give each pair an
index in $\{1,\ldots,L\}$. Every even-multiplicity tuple is generated at
least once, so overcounting gives an upper bound. It follows that
\begin{equation}\label{eq84:moment}
 \limsup_{N\to\infty}\frac1N\sum_{m=0}^{N-1}
 \left|\frac1L\sum_{j=1}^L\mu(m+j)v_j\right|^{2r}
 \leq\frac{(2r-1)!!}{L^r}.
\end{equation}
There are only finitely many terms for fixed $L,r$, so no uniformity in
an increasing number of shifts has been assumed in taking this limit.
The index $m=0$ changes the average by $O(1/N)$ only.

\paragraph{Attempt: choosing the moment order without exchanging limits.}
For a finite net of size $M$, the maximum of the $2r$th powers is at most
the sum of them. H\"older's inequality in $m$, followed by
\eqref{eq84:moment}, turns \eqref{eq84:cover} into
\begin{equation}\label{eq84:tradeoff}
 \limsup_{N\to\infty}\left|\frac1N\sum_{n=1}^N
                      \mu(n)f(T^nx)\right|
 \leq\varepsilon+
 M_{L,\varepsilon}^{1/(2r)}
       \frac{((2r-1)!!)^{1/(2r)}}{\sqrt L}
 \leq\varepsilon+
 \exp\left(\frac{\log M_{L,\varepsilon}}{2r}\right)
                      \sqrt{\frac{2r}{L}}.
\end{equation}
For fixed $\varepsilon$, take
$r=\lceil\sqrt{L\max(1,\log M_{L,\varepsilon})}\rceil$.
By \eqref{eq84:entropy}, $r/L\to0$ and
$\log M_{L,\varepsilon}/r\to0$. The final term in
\eqref{eq84:tradeoff} tends to zero. Thus the limsup is at most
$\varepsilon$, and then $\varepsilon\downarrow0$ proves the target
under \eqref{eq84:chowla}.

The order of operations is essential. For each selected finite $L$ the
integer $r$ and finite net are fixed before $N\to\infty$. Only after that
limit do we increase $L$. Consequently the proof does not assume a
Chowla error term uniform in $r$ or $L$. The zero-entropy cover uses all
points, so the resulting assertion holds for every $x$, not merely almost
every point of an invariant measure. Invertibility of $T$ was never used.

\paragraph{Stress test: what weaker data fail to provide.}
With $r=1$, the same argument gives an error at best
$\sqrt{M_{L,\varepsilon}/L}$, which need not vanish. For instance
$M_L=\exp(\sqrt L)$ is consistent with the subexponential bound
\eqref{eq84:entropy}. Thus zero entropy and a second moment do not by
this argument suffice; the freedom to use increasing finite moment orders
is doing real work. A potentially weaker route than full
\eqref{eq84:chowla} would prove \eqref{eq84:moment} only for the orbit
nets and chosen orders required by \eqref{eq84:tradeoff}, with a similarly
controlled constant. No such uniform arithmetic-net estimate is obtained.

The source about arbitrary prescriptions along a sparse set does not give
a counterexample. If a bounded observable is chosen to correlate perfectly
on $S\subset\N$ of density zero, that part contributes at most
$|S\cap[1,N]|/N\to0$. Positive-density correlation would require additional
structure not supplied by sparse freedom. Likewise logarithmic cancellation
in \eqref{eq84:chowla} would only support a suitably logarithmic version
of the argument; it is not ordinary cancellation. The exact Liouville
linear-form assertion in Section~77 is not silently identified with the
mixed-power M\"obius hypothesis without a separate transfer theorem.

The companion script verifies the combinatorial moment bound by enumerating
independent sign sequences for small $L,r$. This checks the pairing count;
it does not test the required arithmetic cancellation of shifted M\"obius
values at all scales.

\paragraph{Outcome.}
\textbf{Outcome: Conditional reduction.}

An explicit covering/high-moment proof derives every-point ordinary
M\"obius disjointness from the stated mixed-power correlation hypothesis.
It records the entropy--moment tradeoff and the order of limits, and
explains why second moments or sparse-set prescriptions alone do not close
the problem. This reconstructs a classical conditional implication rather
than claiming a new resolution or novelty for that implication.

The missing step is \eqref{eq84:chowla}, or the weaker net-specific
high-moment input just isolated, for the actual arithmetic sequence.
No new unconditional zero-entropy class or counterexample to the original
Sarnak assertion is produced.

% Keep the local bibliography together rather than leave a source fragment.
\needspace{170mm}
% END RESEARCH_ADDITION ATTEMPT_84
\subsection{Sources}
\begin{itemize}\raggedright
\item[\textnormal{[S1]}]\hypertarget{source:84:1}{} Q. Liu, J. Ma, and H. Wang, The Mobius Disjointness Conjecture on infinite-dimensional torus, Introduction (2026).
\url{https://arxiv.org/abs/2603.11087}.
{\small\textit{Catalogue formulation source.}}
\item[\textnormal{[S2]}]\hypertarget{source:84:2}{} Minimal zero entropy subshifts can be unrestricted along any sparse set.
\url{https://www.cambridge.org/core/journals/ergodic-theory-and-dynamical-systems/article/minimal-zero-entropy-subshifts-can-be-unrestricted-along-any-sparse-set/1357B8C6E5361EC6439CEA52634BD7A0}.
{\small\textit{Catalogue research/status reference; not a proof endorsement.}}
% BEGIN RESEARCH_ADDITION REFERENCES_84
\item[\textnormal{[E1]}]\hypertarget{extra:84:1}{} Nikos Frantzikinakis
and Bernard Host, \emph{The logarithmic Sarnak conjecture for ergodic weights},
Annals of Mathematics 187 (2018), 869--931, main theorem and averaging
conventions.
\url{https://annals.math.princeton.edu/2018/187-3/p06}.
{\small\textit{Inspected 27 September 2026; the logarithmic weighting and
system hypotheses are not discarded.}}
\item[\textnormal{[E2]}]\hypertarget{extra:84:2}{} Terence Tao,
\emph{Equivalence of the logarithmically averaged Chowla and Sarnak
conjectures}, arXiv:1605.04628, introduction and comparison with the
ordinary Chowla-to-Sarnak implication.
\url{https://arxiv.org/abs/1605.04628}.
{\small\textit{The ordinary mixed-power implication used here is proved
in full in the notebook; this reference is not treated as a proof of its
unresolved ordinary arithmetic hypothesis.}}
% END RESEARCH_ADDITION REFERENCES_84
\end{itemize}

\end{document}
